%
%Academic CV LaTeX Template (filled for Shuibai Zhang)
%

\documentclass[a4paper,11pt]{article}

% Package imports
\usepackage{latexsym}
\usepackage{xcolor}
\usepackage{float}
\usepackage{ragged2e}
\usepackage[empty]{fullpage}
\usepackage{wrapfig}
\usepackage{lipsum}
\usepackage{tabularx}
\usepackage{titlesec}
\usepackage{geometry}
\usepackage{marvosym}
\usepackage{verbatim}
\usepackage{enumitem}
\usepackage{fancyhdr}
\usepackage{multicol}
\usepackage{graphicx}
\usepackage{cfr-lm}
\usepackage[T1]{fontenc}
\usepackage{fontawesome5}

% Color definitions
\definecolor{darkblue}{RGB}{0,0,139}

% Page layout
\setlength{\multicolsep}{0pt} 
\pagestyle{fancy}
\fancyhf{} % clear all header and footer fields
\fancyfoot{}
\renewcommand{\headrulewidth}{0pt}
\renewcommand{\footrulewidth}{0pt}
\geometry{left=1.4cm, top=0.8cm, right=1.2cm, bottom=1cm}
\setlength{\footskip}{5pt} % Addressing fancyhdr warning

% Hyperlink setup (moved after fancyhdr to address warning)
\usepackage[hidelinks]{hyperref}
\hypersetup{
    colorlinks=true,
    linkcolor=darkblue,
    filecolor=darkblue,
    urlcolor=darkblue,
}

% Custom box settings
\usepackage[most]{tcolorbox}
\tcbset{
    frame code={},
    center title,
    left=0pt,
    right=0pt,
    top=0pt,
    bottom=0pt,
    colback=gray!20,
    colframe=white,
    width=\dimexpr\textwidth\relax,
    enlarge left by=-2mm,
    boxsep=4pt,
    arc=0pt,outer arc=0pt,
}

% URL style
\urlstyle{same}

% Text alignment
\raggedright
\setlength{\tabcolsep}{0in}

% Section formatting
\titleformat{\section}{
  \vspace{-4pt}\scshape\raggedright\large
}{}{0em}{}[\color{black}\titlerule \vspace{-7pt}]

% Custom commands
\newcommand{\resumeItem}[2]{
  \item{
    \textbf{#1}{\hspace{0.5mm}#2 \vspace{-0.5mm}}
  }
}

\newcommand{\resumePOR}[3]{
\vspace{0.5mm}\item
    \begin{tabular*}{0.97\textwidth}[t]{l@{\extracolsep{\fill}}r}
        \textbf{#1}\hspace{0.3mm}#2 & \textit{\small{#3}} 
    \end{tabular*}
    \vspace{-2mm}
}

\newcommand{\resumeSubheading}[4]{
\vspace{0.5mm}\item
    \begin{tabular*}{0.98\textwidth}[t]{l@{\extracolsep{\fill}}r}
        \textbf{#1} & \textit{\footnotesize{#4}} \\
        \textit{\footnotesize{#3}} &  \footnotesize{#2}\\
    \end{tabular*}
    \vspace{-2.4mm}
}

\newcommand{\resumeProject}[4]{
\vspace{0.5mm}\item
    \begin{tabular*}{0.98\textwidth}[t]{l@{\extracolsep{\fill}}r}
        \textbf{#1} & \textit{\footnotesize{#3}} \\
        \footnotesize{\textit{#2}} & \footnotesize{#4}
    \end{tabular*}
    \vspace{-2.4mm}
}

\newcommand{\resumeSubItem}[2]{\resumeItem{#1}{#2}\vspace{-4pt}}

\renewcommand{\labelitemi}{$\vcenter{\hbox{\tiny$\bullet$}}$}
\renewcommand{\labelitemii}{$\vcenter{\hbox{\tiny$\circ$}}$}

\newcommand{\resumeSubHeadingListStart}{\begin{itemize}[leftmargin=*,labelsep=1mm]}
\newcommand{\resumeHeadingSkillStart}{\begin{itemize}[leftmargin=*,itemsep=1.7mm, rightmargin=2ex]}
\newcommand{\resumeItemListStart}{\begin{itemize}[leftmargin=*,labelsep=1mm,itemsep=0.5mm]}

\newcommand{\resumeSubHeadingListEnd}{\end{itemize}\vspace{2mm}}
\newcommand{\resumeHeadingSkillEnd}{\end{itemize}\vspace{-2mm}}
\newcommand{\resumeItemListEnd}{\end{itemize}\vspace{-2mm}}
\newcommand{\cvsection}[1]{%
\vspace{2mm}
\begin{tcolorbox}
    \textbf{\large #1}
\end{tcolorbox}
    \vspace{-4mm}
}

\newcolumntype{L}{>{\raggedright\arraybackslash}X}%
\newcolumntype{R}{>{\raggedleft\arraybackslash}X}%
\newcolumntype{C}{>{\centering\arraybackslash}X}%

% Commands for icon sizing and positioning
\newcommand{\socialicon}[1]{\raisebox{-0.05em}{\resizebox{!}{1em}{#1}}}
\newcommand{\ieeeicon}[1]{\raisebox{-0.3em}{\resizebox{!}{1.3em}{#1}}}

% Font options
\newcommand{\headerfonti}{\fontfamily{phv}\selectfont} % Helvetica-like (similar to Arial/Calibri)
\newcommand{\headerfontii}{\fontfamily{ptm}\selectfont} % Times-like (similar to Times New Roman)
\newcommand{\headerfontiii}{\fontfamily{ppl}\selectfont} % Palatino (elegant serif)
\newcommand{\headerfontiv}{\fontfamily{pbk}\selectfont} % Bookman (readable serif)
\newcommand{\headerfontv}{\fontfamily{pag}\selectfont} % Avant Garde-like (similar to Trebuchet MS)
\newcommand{\headerfontvi}{\fontfamily{cmss}\selectfont} % Computer Modern Sans Serif
\newcommand{\headerfontvii}{\fontfamily{qhv}\selectfont} % Quasi-Helvetica (another Arial/Calibri alternative)
\newcommand{\headerfontviii}{\fontfamily{qpl}\selectfont} % Quasi-Palatino (another elegant serif option)
\newcommand{\headerfontix}{\fontfamily{qtm}\selectfont} % Quasi-Times (another Times New Roman alternative)
\newcommand{\headerfontx}{\fontfamily{bch}\selectfont} % Charter (clean serif font)

\begin{document}
\headerfontiii

%==========================
% Header
%==========================
\begin{center}
    {\Huge\textbf{Shuibai Zhang}}
\end{center}
\vspace{-6mm}

\begin{center}
    \small{
    \href{mailto:shuibai@cs.wisc.edu}{shuibai@cs.wisc.edu} \;|\;
    \href{https://zhangshuibai.github.io}{Personal Website} \;|\;
    \href{https://zhangshuibai.github.io/blog/}{Blog} \;|\;
    \href{https://scholar.google.com/citations?user=jkXzD7YAAAAJ&hl=en&authuser=2}{Google Scholar}
    }
\end{center}
\vspace{-4mm}

\begin{center}
    \small{Madison, Wisconsin, USA}
\end{center}

\vspace{-4mm}

%==========================
% Objective
%==========================
% \section{\textbf{Objective}}
% \vspace{1mm}
% \small{
% Ph.D. student in Computer Science at the University of Wisconsin--Madison, focusing on post-training of large foundation models. I aim to develop efficient, robust, and controllable language models through reinforcement learning with verifiable rewards and diffusion language models for next-generation reasoning and planning systems.
% }
% \vspace{-2mm}

%==========================
% Research Interests
%==========================
\section{\textbf{Research Interests}}
\vspace{1mm}
\small{
I am broadly interested in \textbf{language modeling}, and I explore directions including \textbf{post-training} (e.g., reinforcement learning and on-policy distillation), \textbf{diffusion language models (DLMs)}, and \textbf{test-time training}. 

\vspace{0.5em}
As an impact-driven researcher, I aim to build scalable methods that are effective in practice and deliver real-world impact.
}
\vspace{-2mm}

%==========================
% Education
%==========================
\section{\textbf{Education}}
\vspace{-0.4mm}
\resumeSubHeadingListStart

\resumeSubheading
    {University of Wisconsin--Madison}{Madison, USA}
    {Ph.D. Student in Computer Science}{2024 -- Present}
\resumeItemListStart
    \item Research focus: RL for LLMs, diffusion language models (DLMs), LLM reasoning, Mixtrure of Experts (MOE).
\resumeItemListEnd


\resumeSubheading
{University of Electronic Science and Technology of China (UESTC)}{Chengdu, China}
{B.Eng. in Electronic Information Engineering}{2019 -- 2023}
\resumeItemListStart
\item Dual-degree program with University of Glasgow
\resumeItemListEnd

\resumeSubheading
{University of Glasgow (UoG)}{Glasgow, United Kingdom}
{B.Eng. in Electrical \& Electronic Engineering}{2019 -- 2023}
\resumeItemListStart
\item Dual-degree program with UESTC
\resumeItemListEnd

\resumeSubHeadingListEnd
\vspace{-6mm}

\section{\textbf{Research Highlights}}
\vspace{-0.5mm}
\small{
\begin{itemize}[leftmargin=*,itemsep=3pt]
    \item \textbf{12 accepted papers} (10 at top ML conferences: ICML, ICLR, NeurIPS, COLM, CVPR, ACL), including \textbf{4 oral presentations}; 3 additional manuscripts under review.
    \item \textbf{First or co-first author on 9 papers}, spanning diffusion LLMs, RL post-training, LLM reasoning, and large-scale evaluation.
    \item Experience training 0.5B--7B models; strong background in \textbf{Reinforcement Learning} for LLMs.
    \item Hands-on experience with large-scale LLM training using \textbf{Megatron}, \textbf{VeOmni}, and RL post-training with the \textbf{VERL} framework.
    \item Hands-on experience training and optimizing \textbf{Mixture-of-Experts (MoE)} models, including routing, load balancing, and scaling efficiency.
    \item Extensive experience developing and analyzing \textbf{diffusion LLMs}, covering pretraining, post-training, and inference acceleration.
    \item Reviewer for \textbf{ICLR 2026}, \textbf{ICML 2026}, and \textbf{ACL 2026}.
\end{itemize}
}




\section{\textbf{Open-Source Contributions}}
\vspace{-0.5mm}
\resumeSubHeadingListStart

\resumeProject
  {Open-dCoder: The First Fully Open Diffusion LLM for Code}
  {Open Source Project}
  {Jul 2025 -- Sep 2025}
  {{[\href{https://github.com/pengzhangzhi/Open-dLLM}{\textcolor{darkblue}{GitHub (\faStar{} 632)}}] \; [\href{https://oval-shell-31c.notion.site/Open-dLLM-Open-Diffusion-Large-Language-Model-25e03bf6136480b7a4ebe3d53be9f68a}{\textcolor{darkblue}{Project Page}}]}}
\resumeItemListStart
  \item \textbf{Co-Lead}: Co-directed the development and release of the first fully open-source diffusion LLM for coding.
  \item Built and released a 0.5B diffusion language model for code with full pipeline and evaluation suite.
  \item Project has received strong community adoption with \textbf{579 GitHub stars} by Apr 2026.
\resumeItemListEnd
\resumeSubHeadingListEnd

% \section{\textbf{Research Projects}}
% \vspace{-0.4mm}
% \resumeSubHeadingListStart

% ===================== Project 1 =====================
% \resumeProject
%   {Diffusion Language Model Beats Autoregressive Language Model in Correction}
%   {Sole Lead Researcher — Ongoing}
%   {Sep 2025 -- Present}
%   {}
% \resumeItemListStart
%   \item \textbf{Sole lead researcher} overseeing problem formulation, benchmark design, model development, and experiments.
%   \item Developed the \textbf{Code Revision Benchmark (CRB)}, a new evaluation suite for non-causal ``correction'' tasks such as bug fixing and constraint-based puzzle resolution.
%   \item Introduced a new \textbf{post-training paradigm for diffusion LLMs} that enables iterative error correction by conditioning on corrupted inputs, and developed a \textbf{correction-oriented diffusion language model (cDLM)} built on this framework.
%   \item Preparing manuscript for submission to \textbf{ICML 2026}.
% \resumeItemListEnd



% ===================== Project 2 =====================
% \resumeProject
%   {Is the Importance Ratio Necessary for Stable Reinforcement Learning in LLMs?}
%   {Sole Lead Researcher — Krafton AI}
%   {Jul 2025 -- Sep 2025}
%   {{[\href{https://drive.google.com/file/d/1ci-g7yh1u5dR_jGvNCZAAoUk84fMulOw/view?usp=drive_link}{\textcolor{darkblue}{Manuscript}}]}}
% \resumeItemListStart
%   \item \textbf{Sole lead researcher} driving theory, algorithm design, and experimental validation.
%   \item Identified that clipping—not importance-ratio scaling—is the primary stabilizer in Group Relative Policy Optimization (GRPO) for LLM post-training.
%   \item Proposed \textbf{Self-Importance Clipping (SIC)}, a \textbf{ratio-free} optimization method leveraging model confidence as a gating signal to suppress unstable updates.
%   \item Provided formal analysis showing that extreme-confidence tokens amplify KL sensitivity, explaining how clipping enforces a soft trust-region constraint on policy drift.
%   \item Achieved \textbf{37\% faster training} while maintaining GRPO-level stability and reasoning performance on AIME, OlympiadBench, and MATH500 using Qwen models.
% \resumeItemListEnd

% \resumeSubHeadingListEnd
% \vspace{-6mm}


\section{\textbf{Publications \& Manuscripts}}
\vspace{0.5mm}
\small{\textit{* = equal contribution}}
\vspace{1mm}

\small{
\begin{itemize}[leftmargin=*,itemsep=4pt]

% ===== Manuscripts Under Review (newest first) =====
\item \href{https://arxiv.org/abs/2608.18597}{
\textbf{Shuibai Zhang}*, Xinchi Liu*, Fred Zhangzhi Peng, Zhihan Yang, Shutong Wu, Yingzi Ma, Jiawei Zhang,
``Off-Manifold Collapse in Guided Protein Language Models''} \\
\hfill \textit{Under Review}

\item \href{https://arxiv.org/abs/2604.11978}{
Xinyu Jessica Wang, Haoyue Bai, Yiyou Sun, Haorui Wang, \textbf{Shuibai Zhang}, Wenjie Hu, Mya Schroder, Bilge Mutlu, Dawn Song, Robert D. Nowak,
``The Long-Horizon Task Mirage? Diagnosing Where and Why Agentic Systems Break''} \\
\hfill \textit{Under Review --- EMNLP 2026}

\item \href{https://arxiv.org/pdf/2512.15596}{
\textbf{Shuibai Zhang}, Fred Zhangzhi Peng, Yiheng Zhang, Jin Pan, Grigorios G. Chrysos,
``Corrective Diffusion Language Model''} \\
\hfill \textit{Under Review --- NeurIPS 2026}

% ===== Accepted Publications (newest first) =====

\item \href{https://arxiv.org/pdf/2605.18530}{
Zhihan Yang, Wei Guo, \textbf{Shuibai Zhang}, Subham Sekhar Sahoo, Yongxin Chen, Arash Vahdat, Morteza Mardani, John Thickstun,
``Continuous Diffusion Scales Competitively with Discrete Diffusion for Language''} \\
\hfill \textit{COLM 2026 Workshop on Non-Autoregressive Language Models --- \textbf{Oral Presentation}}

\item \href{https://arxiv.org/abs/2604.01622}{
\textbf{Shuibai Zhang}*, Caspian Zhuang*, Chihan Cui*, Zhihan Yang, Fred Zhangzhi Peng, Yanxin Zhang, Haoyue Bai, Zack Jia, Yang Zhou, Guanhua Chen, Ming Liu,
``Expert-Choice Routing Enables Adaptive Computation in Diffusion Language Models''} \\
\hfill \textit{COLM 2026}

\item \href{https://www.arxiv.org/pdf/2512.00831}{
Yuchen Zeng*, \textbf{Shuibai Zhang}*, Wonjun Kang*, et al., Dimitris Papailiopoulos, Kangwook Lee,
``ReJump: A Tree-Jump Representation for Analyzing and Improving LLM Reasoning''} \\
\hfill \textit{ICML 2026}

\item \href{https://openreview.net/forum?id=eiOwyAa2O1}{
\textbf{Shuibai Zhang}*, Junhyuck Kim*, Gyeongman Kim, Jaewoong Cho,
``Is the Importance Ratio Necessary for Stable Reinforcement Learning in LLMs?''} \\
\hfill \textit{ICLR 2026 SPOT Workshop: Scaling Post-training for LLMs}

\item \href{https://www.arxiv.org/pdf/2509.23405}{
Fred Zhangzhi Peng, Zachary Bezemek, Jarrid Rector-Brooks, \textbf{Shuibai Zhang}, Anru R. Zhang, Michael Bronstein, Avishek Joey Bose, Alexander Tong,
``Planner Aware Path Learning in Diffusion Language Models Training''} \\
\hfill \textit{ICLR 2026 --- \textbf{Oral Presentation}}

\item \href{https://arxiv.org/pdf/2510.04767}{
Wonjun Kang, Kevin Galim, Seunghyuk Oh, Minjae Lee, Yuchen Zeng, \textbf{Shuibai Zhang}, Coleman Hooper, Yuezhou Hu, Hyung Il Koo, Nam Ik Cho, Kangwook Lee,
``ParallelBench: Understanding the Tradeoffs of Parallel Decoding in Diffusion LLMs''} \\
\hfill \textit{ICLR 2026}

\item \href{https://arxiv.org/pdf/2512.04459v1}{
Yingzi Ma, Yulong Cao, Wenhao Ding, \textbf{Shuibai Zhang}, Yan Wang, Boris Ivanovic, Ming Jiang, Marco Pavone, Chaowei Xiao,
``dVLM-AD: Enhance Diffusion Vision-Language-Model for Driving via Controllable Reasoning''} \\
\hfill \textit{CVPR 2026 (Findings)}

\item \href{https://arxiv.org/pdf/2502.06737}{
Thomas Zeng, \textbf{Shuibai Zhang}, et al., Kannan Ramchandran, Dimitris Papailiopoulos, Kangwook Lee,
``VersaPRM: Multi-Domain Process Reward Model via Synthetic Reasoning Data''} \\
\hfill \textit{ICML 2025 --- \textbf{Oral Presentation}}

\item \href{https://arxiv.org/pdf/2410.08695}{
Yue Yang*, \textbf{Shuibai Zhang}*, Wenqi Shao*, Kaipeng Zhang, Yi Bin, Yu Wang, Ping Luo,
``Dynamic Multimodal Evaluation with Flexible Complexity by Vision-Language Bootstrapping''} \\
\hfill \textit{ICLR 2025 --- \textbf{Oral Presentation}}

\item \href{https://arxiv.org/pdf/2406.07230}{
Weiyun Wang*, \textbf{Shuibo Zhang}*, Yiming Ren*, Yuchen Duan*, Tiantong Li*, et al., Wenhai Wang,
``Needle In A Multimodal Haystack''} \\
\hfill \textit{NeurIPS 2024}

\item \href{https://arxiv.org/pdf/2312.15918}{
Linyi Yang*, \textbf{Shuibai Zhang}*, Zhuohao Yu*, et al., Xing Xie, Weizhu Chen, Yue Zhang,
``Supervised Knowledge Makes Large Language Models Better In-Context Learners''} \\
\hfill \textit{ICLR 2024}

\item \href{https://arxiv.org/pdf/2211.08073}{
Linyi Yang*, \textbf{Shuibai Zhang}*, et al., Jindong Wang, Xing Xie, Yue Zhang,
``GLUE-X: Evaluating Natural Language Understanding Models from an Out-of-distribution Generalization Perspective''} \\
\hfill \textit{ACL 2023 (Findings)}

\end{itemize}
}
\vspace{-4mm}



%==========================
% Work Experience
%==========================
\section{\textbf{Research Experience}}
\vspace{-0.4mm}
\resumeSubHeadingListStart

\resumeSubheading
{\href{https://www.krafton.ai/en/}{Krafton AI}}{Madison, USA}
{Research Intern}{Jul 2025 -- Sep 2025}

\resumeSubheading
{\href{https://www.shlab.org.cn/}{Shanghai Artificial Intelligence Laboratory}}{Shanghai, China}
{Research Assistant}{May 2024 -- Aug 2024}

\resumeSubheading
{\href{https://westlakenlp.com/}{WestlakeNLP Lab, Westlake University}}{Hangzhou, China}
{Research Assistant}{Jun 2023 -- Apr 2024}

\resumeSubHeadingListEnd
\vspace{-6mm}

%==========================
% Teaching Experience
%==========================
\section{\textbf{Teaching Experience}}
\vspace{-0.4mm}
\resumeSubHeadingListStart

\resumeSubheading
{COMP SCI 220: Data Science Programming I}{University of Wisconsin--Madison}
{Teaching Assistant}{Spring \& Fall 2025}

\resumeSubheading
{COMP SCI 300: Programming II}{University of Wisconsin--Madison}
{Teaching Assistant}{Fall 2024}

\resumeSubHeadingListEnd
\vspace{-4mm}

\section{\textbf{Technical Skills}}
\vspace{-0.5mm}
\small{
\begin{itemize}[leftmargin=*,itemsep=3pt]

    \item \textbf{Language Models:} Diffusion LLMs, Autoregressive LLMs, RLHF/RLVR, GRPO/PPO/SFT, MOE Design.

    \item \textbf{Machine Learning:} Deep Learning, Reinforcement Learning, Optimization, Evaluation Design.

    \item \textbf{Large-Scale Training:} Distributed training (DDP/FSDP); proficient with \textbf{Megatron} and the \textbf{VeOmni} pretraining framework, and the \textbf{VERL} RL post-training framework.

    \item \textbf{Tools \& Infrastructure:} Weights \& Biases, Docker, Git, Linux.

\end{itemize}
}
\vspace{-4mm}



\end{document}
